\section{Complexity and Scalability Analysis} \label{sec:complexity-and-scalability-analysis}
\subsection{Variable Growth} \label{subsec:variable-growth}

Let:
\begin{align*}
    n &= \text{number of floors}, \\
    m &= \text{number of movers}, \\
    k &= \text{number of items}, \\
    T &= \text{time horizon}.
\end{align*}

\noindent For Task 1, the number of variables is approximately:
\[
    m n (T+1) + 2mT.
\]
The first term represents mover positions, while the second term represents the ascend and descend action variables.

\noindent For Task 2, item positions and carrying variables are added:
\[
    m n (T+1) + k n (T+1) + 2mT + mkT.
\]

\noindent Thus, the number of variables grows linearly with the time horizon $T$, but also depends on the product $mk$ because every mover-item pair may correspond to a carrying variable at every step.

\subsection{Clause Growth} \label{subsec:clause-growth}

The largest source of clauses comes from exactly-one and at-most-one constraints.
Because the implementation uses pairwise at-most-one clauses, enforcing uniqueness over $r$ literals requires:
\[
    \binom{r}{2}
\]
binary clauses.

\noindent For example, requiring a person to be on exactly one of $n$ floors at a given time step requires:
\[
    1 + \binom{n}{2}
\]
clauses.
This is repeated for every person and every time step.

\noindent Task 2 adds similar constraints for items, as well as carrying mutex constraints between movers and items.
Therefore, the number of clauses grows polynomially in the input parameters for a fixed horizon.

\subsection{SAT Solving Complexity} \label{subsec:sat-solving-complexity}

\noindent Although the encoding size is polynomial in $n$, $m$, $k$, and $T$, SAT solving is NP-complete in the worst case.
Therefore, there is no guarantee that every large instance will be solved quickly.

\noindent In practice, modern SAT solvers such as Glucose3 exploit structure, unit propagation, clause learning, and branching heuristics.
For the tested problem sizes, the solver usually returns results within milliseconds.

\subsection{Practical Scalability} \label{subsec:practical-scalability}

The most important factor affecting performance is the chosen time horizon.
In fixed-horizon mode, an unrealistically small $T$ often produces an unsatisfiable instance quickly, while a larger $T$ may produce many more variables and clauses.
In iterative deepening mode, the backend solves a sequence of formulas, which is useful for optimality but can become expensive when the optimal plan requires many steps.